\scanpage{13}

$M(S)$ définissent des courbes elliptiques isomorphes (non structurées) il faut que \deleted{\unsure{donc}} ils soient conjugués par une opération de $G(S)$, où \inserted{$G$ est un groupe \unsure{linéaire|lisse} sur $S$.} Enfin, pour tte courbe elliptique $C$ sur $S$, il existe une extension \deleted{fid. plat}\inserted{simple surj. de type fini} $S'$ de $S$ [savoir le fibré principal associé à $\mathcal E_3$ pour le \unsure{clivage} qu'on a mis sur $\mathcal E_3$] telle que \deleted{sur} $C_{S'}$ ``provienne'' d'un élément de $M(S')$.

Ainsi, avec une terminologie générale, le foncteur $F$ : ``classes à isom près de courbes elliptiques sur $S$'' est \deleted{\unsure{pro}}-$\mathcal D$-représentable, où $\mathcal D$ désigne la catégorie des morphismes simples surjectifs. Si ce n'était que pour \unsure{obtenir cet énoncé}, on n'aurait pas à se donner tant de mal, \uline{car} on a un raisonnement général qui le donne pour les courbes de genre $g$ \unsure{quelc} (de façon \unsure{triangulaire} \unsure{quotient} par un groupe projectif---).
\mspage{12}
